\documentclass[a4paper]{article}

\usepackage[spanish]{babel}
\usepackage{graphicx}
\usepackage{amsmath, amssymb}
\usepackage[margin=2cm]{geometry}
\usepackage{fancyhdr}
\usepackage{enumerate}
\usepackage[shortlabels]{enumitem}
\usepackage{parskip}
\usepackage[most]{tcolorbox}
\usepackage[hidelinks]{hyperref}
\usepackage{float}
\usepackage{booktabs}
\usepackage{mathrsfs}
\usepackage{tikz}

% cabecera
\pagestyle{fancy}
\fancyhead[l]{Mecánica celeste}
\fancyhead[c]{Quiz 2}
\fancyhead[r]{\today}
\fancyfoot[c]{\thepage}
\renewcommand{\headrulewidth}{0.2pt} % linea horizontal


\begin{document}


\section{Problema 1: Energía y sistema ligado}
Por qué razón $E<0$ no necesariamente implica que se tiene un sistema ligado.

\section{Problema 2: Términos de la energía potencial}

Demuestre que el número de términos necesarios para calcular la energía potencial de un sistema de $N$ cuerpos es $\frac{N(N-1)}{2}$ y por lo tanto su costo computacional es $\mathcal{O}(N^2)$

\section{Problema 3: Dibujar árbol jerárquico}

Dibuje el árbol jerárquico del sistema sextuple TYC 7037-89-1

\section{Problema 4: Conservación de la energía total de un sistema}

Partiendo de $$\left\{m_i \ddot{\vec{r}}_i = -\frac{\partial U}{\partial \vec{r}_i}\right\}_N$$ demuestre que la energía total del sistema $E=K+U$ se conserva

\section{Problema 5: Demostración analítica}

Demuestre que si $G = \sum m_i \vec{r}_i\cdot \dot{\vec{r}}_i$ entonces $\dot{G} = K+E$ 

\textit{Ayuda:} $U(x_1,y_1,z_1,x_2,...,z_N)$ es una función homogénea de orden $k=-1$

\section{Problema 6: Coordenadas}
Partiendo de las definiciones del vector relativo $\vec{r}=\vec{r}_1-\vec{r}_2$ y del centro de masa $\vec{R}_{cm}$ demuestre que 

\begin{align*}
\vec{r}_1 &= \vec{R}_{cm} + \frac{m_2}{M}\vec{r} \\
\vec{r}_2 &= \vec{R}_{cm} - \frac{m_1}{M}\vec{r}
\end{align*}

\bibliographystyle{plainurl}
\bibliography{bibliografia} 
\end{document}
